\section{Мультиобъектное отслеживание}\label{sec::mot}

Мультиобъектное отслеживание (англ. Multi Object Tracking (MOT)) — это процесс идентификации и отслеживания нескольких объектов в пространстве и времени, при помощи компьютерного зрения. MOT важен для понимания и анализа сложных сцен, где множество объектов взаимодействуют друг с другом.

В ИСПВ МОТ играет ключевую роль, так как его задача обеспечивать устойчивое определение положения, скорости и идентичности объектов с течением времени. Эта информация критически необходима для таких ключевых функций ИСПВ, как автоматическое экстренное торможение (АЭТ), адаптивный круиз-контроль (АКК), система удержание в полосе (СУП) и т.д.

Мультиобъектное отслеживание выполняет функцию синхронизации и интеграции данных от различных сенсоров, обеспечивая временную и пространственную связность между измерениями. Это позволяет формировать целостную модель окружающей среды, что, в свою очередь, служит основой для построения траекторий и выработки управляющих воздействий на исполнительные органы транспортного средства.

МОТ может быть применим ко всем уровням автономности, кроме нулевого, в той или иной степени. Например, для уровней 1-2 MOT может использоваться для таких функций, как АКК и АЭТ, а уже для уровней с третьего и выше, MOT является критически важной составляющей, обеспечивающей осведомленность систем автомобиля о динамической дорожной обстановке.

\subsection{Основные сенсоры, применяемые в МОТ}\label{subsec:sensors}

Для повышения безопасности и эффективности движения, система помощи водителю должна иметь точную и подробную информацию об окружающей среде, которую она получает с датчиков. В зависимости от применяемого типа датчика, система получает определенную информацию о мире. В данной части будут рассмотрены основные сенсоры, которые применяются в ИСПВ, их преимущества и недостатки, а также возможность комбинирования друг с другом.

\subsubsection{Камеры}

Камера -- один из самых важных сенсоров для получения информации об окружающем мире, так как они формируют высокодетализированное цветовое изображение, пригодное для распознавания дорожной разметки, знаков и классификации объектов. Слабым местом камеры является чувствительность к погодным условиям и освещенности сцены.
 
\subsubsection{Радары}

Радар -- в рамках ИСПВ, данный сенсор позволяет выделять объекты вокруг ТС, измерять расстояние до них и скорость. Так как работа происходит в радиодиапазоне, радары практически не чувствительны к погодным условиям, однако имеют низкое угловое разрешение (градусы вместо угловых минут) и трудности с классификацией целей из-за слабого сигнального контраста.

\subsubsection{Лидары}

Лидар -- один из самых универсальных датчиков, так как он предоставляет довольно 3-D геометрию окружающей среды и мало зависит от освещённости. Однако из-за высокой плотности сканирования, работа лидара может быть нарушена плохими погодными условиями. Также, цена этого датчика довольно высокая и его надежность во многом зависит от его конструкции. 
 
\subsubsection{Сенсорное слияние}
Для нивелирования недостатков каждого из сенсоров, используют сенсорное слияние (Sensor Fusion)~\cite{Liu2024}. Данная технология объединяет данные с различных датчиков, тем самым повышая точность и надежность данных, характеризующих окружающую среду. 

Комбинация камеры и радара позволит довольно точно распознавать дорожные знаки, разметку, объекты и при этом иметь информацию о расстоянии и скорости объектов вне зависимости от погодных условий. Такая комбинация является наиболее выгодной в плане стоимости и функциональности. 

Камера и лидар позволяют выстраивать 3D карту окружающей среды с цветовой и текстурной информацией для улучшения классификации. Однако восприимчивость к погодным условиям является серьезным недостатком. 

Радар и лидар дают возможность ИСПВ быть невосприимчивой к погодным условия и освещению, а также создавать 3D карту окружающей среды благодаря высокому угловому разрешению. Но при такой комбинации пропадает возможность определять разметку и дорожные знаки. 

Комбинация всех трех датчиков позволит использовать все доступные данные для наиболее полного восприятие окружающей среды. Несмотря на это, выбор такого типа сенсорного слияния требует больших вычислительных мощностей и сильно увеличивает стоимость системы в целом.

Таким образом, основываясь на опыте коллег программистов и инженеров, можно прийти к выводу, что использование камеры и радара покрывает практически весь спектр задач. Данная комбинация обеспечивает не только высокую точность определения расстояния и скорости объектов, но и позволяет эффективно распознавать все объекты, начиная от едущих рядом автомобилей, заканчивая дорожными знаками и разметкой. Относительно небольшая стоимость также является важным фактором выбора. 

В случае, если требования к точности и надежности очень высокие, возможна и интеграция лидара в систему сенсоров. 

В рамках данной работы система МОТ рассматривается только с использованием камеры.

\subsection{Категоризация МОТ}\label{subsec:MOT_category}

В виду того, что МОТ является довольной комплексной областью исследования, разделение на группы является достаточно сложной задачей. Категоризация на следующие группы является наиболее универсальной, так как охватывает основные черты МОТ~\cite{Article::Luo2021}: 

\begin{enumerate}

	\item Категоризация по методу инициализации: 
		
	\begin{enumerate}
		
		\item Слежение с обнаружением (Detection-Based Tracking (DBT))~\cite{Bose2007}~\cite{Song2010}: Объекты сначала обнаруживаются, а затем связываются в траектории. Этот метод фокусируется на конкретных типах целей, таких как пешеходы, транспортные средства или лица;
		
		\item Слежение без обнаружения (Detection-Free Tracking (DFT))~\cite{Zhang2013}: Требует ручной инициализации фиксированного количества объектов в первом кадре, затем локализует эти объекты в последующих кадрах.
	
	\end{enumerate}
	
	\item Категоризация по режиму обработки:
	
	\begin{enumerate}
		
		\item Онлайн-отслеживание~\cite{Xiang2015}: Обрабатывает последовательность изображений пошагово, используя наблюдения, доступные до текущего момента;
		
		\item Оффлайн-отслеживание: Использует пакет кадров для обработки данных, требует предварительного получения наблюдений со всех кадров.
		
	\end{enumerate}
	
	\item Категоризация по типу вывода:
	
	\begin{enumerate}
		
		\item Стохастическое отслеживание: Результаты варьируются от одного запуска к другому из-за случайности в процессе генерации частиц или других элементов;
		
		\item Детерминированное отслеживание: Результаты остаются постоянными при многократном запуске методов.
		
	\end{enumerate}
	
\end{enumerate}

Так как система МОТ будет использована в автономном транспорте, то нам необходимы алгоритмы слежения с обнаружением в онлайн режиме. Таким образом, надо подобрать такие системы детекции и трекинга, которые будут удовлетворять нашим требованиям.

\subsection{Методы детекции}\label{subsec:detector}

Проблема детекции объекта на изображении или видео, стоит давно. Современные технологии позволяют довольно точно определять что изображено на картинке или видео. Рассмотрим несколько популярных методов детекции и проведем сравнение.

\subsubsection{Метод градиентов (Histogram of Oriented Gradients (HOG))}

Метод градиентов появился в 2005 году Навнеетом Далалом и Биллом Триггсом. Данный способ заключается в том, что с помощью распределения градиентов интенсивности или направлений краев определяется форма объекта и его внешний вид. На изображениях \ref{fig::HOG_1} и \ref{fig::HOG_2} показано как распределяется градиент в изображении.

\img[htb]{fig::HOG_1}{HOG_1.png}{Распределение векторов градиента}{0.8}

Из плюсов можно выделить вычислительную эффективность: метод довольно эффективен и может быть реализован для работы в реальном времени на современном аппаратном обеспечении и имеет широкую область применения для задач обнаружения пешеходов, транспортных средств, лиц и других объектов. 

\img[htb]{fig::HOG_2}{HOG_2.png}{Пример определения формы объекта при помощи HOG}{0.8}

Однако у HOG есть ряд минусов, которые не позволяют его использовать для МОТ: игнорирование цвета, что не позволит различить объекты друг от друга, неустойчивость к изменению масштаба и ориентации объекта, чувствительность к шумам и плохая работа с динамическими сценами.  

\subsubsection{Сверточные нейронные сети (СНС)}

Сверточные нейронные сети --- это класс глубоких нейронных сетей, которые используются в компьютерном зрении для анализа визуальных данных. 

Благодаря своим алгоритмам, СНС способны распознавать объекты независимо от их положения в изображении, а также способны автоматически выявлять важные признаки без предварительного определения их в ручную. Но из-за высокой вычислительной сложности СНС требует значительных вычислительных ресурсов и большего количества времени на обработку изображения, что является критически важным фактором для МОТ. Долгая обработка не позволит использовать данный метод в реальном времени.

\subsubsection{Single Shot Detector (SSD)}

Single Shot Detector (SSD) --- это метод детекции объектов, используемый в компьютерном зрении для обнаружения множества объектов на изображении в один проход. Это делает данный способ особенно полезным для приложений в реальном времени, таких как видеонаблюдение или автономное вождение, где требуется быстрое и эффективное распознавание. Одним из главных преимуществ SSD является скорость. Благодаря своим алгоритмам, SSD может оценить сетку ограничивающих рамок и вероятности классов для этих рамок одновременно, что позволяет достичь скорости, которую можно использовать в реальном времени. Также, благодаря архитектуре, SSD проще в реализации и тренировке по сравнению с вышеперечисленными методами.

Несмотря на все преимущества, у SSD есть и ряд недостатков, таких как сложность работы с мелкими объектами и высокая степень ложных положительных срабатываний (особенно в сценах с сложным фоном).

SSD продолжает развиваться, и его варианты, такие как SSD300 (использующий входные изображения размером 300x300 пикселей) и SSD512 (использующий изображения размером 512x512), предлагают различные компромиссы между скоростью и точностью детекции.

\subsubsection{YOLO(You Only Look Once)}

YOLO или You Only Look Once --- алгоритм детекции объектов, схожий с SSD, также обрабатывает изображение только один раз. Однако, между ними есть ряд различий, таких как различие в архитектуре, скорости и точности. Различие в архитектуре определяет скорость. YOLO делит изображение на сетку, а затем использует набор фильтров, чтобы предсказать все рамки и классификации за один раз (рисунок \ref{fig::YOLO_1}). SSD также предсказывает ограничивающие рамки и классы объектов за один проход, но делает это, применяя серию сверточных фильтров различного размера к множеству признаковых карт, что позволяет обнаруживать объекты разного размера. Таким образом, YOLO более эффективен для обнаружения больших объектов, но может пропустить маленькие из-за фиксированного размера сетки, что не является критичным в рамках автономных транспортных средств. 

\img[htb]{fig::YOLO_1}{YOLO_1.png}{Принцип работы YOLO}{1}

Одним из ключевых факторов это скорость. SSD уступают по скорости YOLO, которая способна обрабатывать от десятков до сотен кадров в секунду. 

Таким образом, после небольшого анализа и сравнения методов детекции, наиболее подходящим  в МОТ для автономных транспортных средств является YOLO, благодаря своей высокой скорости и точности.

\subsection{Основные метрики МОТ}\label{subsec:metrics}

Для дальнейшего понимания, подходят ли те или иные алгоритмы, необходимо определить их преимущества и недостатки, конкретно под требуемую задачу. Для этого надо пользоваться метриками, которые покажут в чем заключается различие подобранных алгоритмов.

\textbf{Полнота (Recall)} --- доля правильно обнаруженных объектов от общего количества истинных объектов:  
  \[
     \text{Recall}= \frac{\mathrm{TP}}{\mathrm{TP}+\mathrm{FN}} .
  \]
  Рост Recall означает, что трекер реже пропускает цели, но не говорит о числе ложных тревог.

\textbf{Точность (Precision)} --- доля верных детекций среди всех выданных алгоритмом:  
  \[
     \text{Precision}= \frac{\mathrm{TP}}{\mathrm{TP}+\mathrm{FP}} .
  \]
  Высокая Precision возможна даже при малом охвате, если трекер «осторожничает» и подаёт мало боксов.

\textbf{FAF (False Alarms per Frame)} --- среднее число ложных срабатываний на кадр:  
  \[
     \text{FAF}= \frac{\mathrm{FP}}{\text{\#\,кадров}} .
  \]
  Удобна в робототехнике, поскольку напрямую показывает, сколько «фантомов» придется обрабатывать downstream-модулю.

\textbf{MODA (Multiple-Object Detection Accuracy)} --- комплексная метрика для детекции без учёта идентичностей:  
  \[
     \text{MODA}=1-\frac{\mathrm{FN}+\mathrm{FP}}{\mathrm{GT}} ,
  \]
  где \(\mathrm{GT}\) — число объектов по разметке. Суммирует пропуски и ложные тревоги, но не различает их происхождение.

\textbf{MOTP (Multiple-Object Tracking Precision)} --- оценивает точность локализации предсказанных боксов относительно эталона:  
  \[
     \text{MOTP}= \frac{\sum_{(i,t)}\text{IoU}_{i,t}}{\text{\#\,соответствий}} ,
  \]
  где \(\text{IoU}_{i,t}\) — перекрытие «истины» и предсказания для объекта \(i\) в кадре \(t\).

\textbf{MOTA (Multiple-Object Tracking Accuracy)} --- самая популярная агрегированная величина:  
  \[
     \text{MOTA}=1-\frac{\mathrm{FN}+\mathrm{FP}+\mathrm{IDSW}}{\mathrm{GT}} .
  \]
  Сводит воедино пропуски, ложные срабатывания и переключения идентификаторов, но «награждает» хороший детектор сильнее, чем хорошую ассоциацию.

\textbf{IDS (ID Switches)} --- число раз, когда активная траектория меняет присвоенный ей идентификатор. Метрика-счётчик, не нормируется и используется для прямого сравнения трекеров на одной последовательности.

\textbf{IDF1} --- F-мера для сохранения идентичности; рассматривает множества пар «кадр–ID»:  
  \[
     \text{IDF1}= \frac{2\,\mathrm{IDTP}}{2\,\mathrm{IDTP}+\mathrm{IDFP}+\mathrm{IDFN}} ,
  \]
  где \(\mathrm{IDTP}\) — кадры, в которых и объект, и его ID предсказаны верно. Слабо зависит от точности размещения боксов.

\textbf{HOTA (Higher-Order Tracking Accuracy)} --- современный компромисс, равновешивающий детекцию и ассоциацию:  
  \[
     \text{HOTA}= \sqrt{\text{DetA}\times\text{AssA}} .
  \]
  Здесь \(\text{DetA}\) оценивает корректность детекций (TP/FP/FN), а \(\text{AssA}\) — корректность сопоставления идентификаторов. Геометрическое среднее гарантирует, что провал в любой из задач тянет итоговое значение вниз.

\subsection{Алгоритмы МОТ}\label{subsec:mot_algotithms}

Разработка надёжной системы мультиобъектного трекинга (MOT) для автономных транспортных средств требует выбора подходящего базового алгоритма, который можно модифицировать и адаптировать к конкретным условиям. На сегодняшний день в научной и прикладной среде существует ряд популярных трекеров, различающихся по архитектуре, качеству отслеживания, устойчивости к заслонам и вычислительной сложности.

В таблице~\ref{tab:mot_comparison} представлено сравнение шести наиболее известных алгоритмов трекинга, получивших широкое распространение в задачах MOT. Все трекеры запускались на едином детекторе (YOLOv5) и тестировались на подмножестве MOT17 с идентичными параметрами~\cite{zhang2021bytetrack, cao2022observation, bewley2016simple, zhang2020fairmot, meinhardt2022trackformer, zeng2022motr}.

\begin{table}[h!]
\centering
\caption{Сравнение трекеров по ключевым метрикам на MOT17}
\label{tab:mot_comparison}
	\begin{tabular}{|p{3.5cm}|p{2cm}|p{2cm}|p{2cm}|p{2cm}|p{2cm}|}
		\hline
		\textbf{Трекер} 	& MOTA ↑ 		& IDF1 ↑ 		& HOTA ↑		& FAF ↓			& FPS ↑			\\
		\hline
		ByteTrack 			& \textbf{80.3} & 77.3			& 63.1 			& \textbf{0.18}	& 30  			\\
		OC-SORT 			& 78.0 			& \textbf{77.5} & 63.2 			& 0.21 			& \textbf{700} 	\\
		FairMOT 			& 67.9 			& 68.1 			& 56.3 			& 0.31 			& 28  			\\
		TrackFormer 		& 74.1 			& 68.0 			& 60.0 			& 0.25 			& 5.7  			\\
		MOTR 				& 71.9 			& 68.4 			& 57.2 			& 0.29			& 6 			\\
		\textbf{DeepSORT} 	& 78.8 			& 76.5 			& \textbf{63.5}	& 0.19			& 80 			\\
	\hline
	\end{tabular}
\end{table}

Исходя из данных, приведенных в таблице~\ref{tab:mot_comparison}, можно выделить три алгоритма, метрики которых превосходят другие: ByteTrack, OC-SORT и DeepSORT. Наиболее мощным можно выделить алгоритм ByteTrack, показатели которого довольно высоки, однако, модель построена на строго геометрической модели движения и не использует признаковую ифнормацию для различия объектов. Такая же ситуация наблюдается и в OC-SORT, в которой также не использует признаковую ифнормацию. Это делает их уязвимыми в условиях визуального сходства и частых заслонов. Кроме того, они ограничены жёсткой схемой линейного движения, что снижает устойчивость при сложной кинематике объектов.

\subsubsection{Преимущества DeepSORT как базовой системы}

В данном проекте в качестве базового трекера остановимся на алгоритме DeepSORT. Он представляет собой расширение классического SORT, в которое встроен модуль ReID (визуальных эмбеддингов) и фильтр Калмана. Это обеспечивает:

\begin{itemize}
    \item устойчивость к идентификационным ошибкам в условиях визуальной близости объектов;
    \item хорошую модульность — возможность заменить детектор или ReID-сеть без изменения остального пайплайна;
    \item прозрачность и интерпретируемость всех шагов (в отличие от end-to-end моделей);
    \item разумный компромисс между точностью и вычислительной сложностью.
\end{itemize}

Однако стандартный DeepSORT использует линейную модель движения, которая плохо справляется с резкими поворотами, ускорением и сменой масштаба объекта. Это создаёт условия для ID-переключений и обрывов треков.

С целью устранения указанных недостатков в рамках данной работы предлагается модифицировать DeepSORT, заменив стандартный линейный фильтр Калмана на расширенный (Extended Kalman Filter, EKF). Такой фильтр способен учитывать нелинейную динамику движения объектов и более точно корректировать состояние трека при получении новых наблюдений.

Ожидается, что внедрение EKF позволит:
\begin{itemize}
    \item улучшить устойчивость к пропаданию объектов за заслонами;
    \item сократить число ID-switch за счёт более точной предсказательной модели;
    \item повысить HOTA и IDF1 без ущерба для скорости работы.
\end{itemize}

Таким образом, выбор DeepSORT как отправной точки для разработки обоснован его архитектурной гибкостью, хорошими исходными показателями и потенциалом для усовершенствования. Представленная в данной работе модификация нацелена на устранение выявленных ограничений и достижение конкурентного качества в условиях городской среды.

